% Appendix B: Linear Probing Training Details

\section{Linear Probing Training Details}
\label{app:probing}

This appendix details the training procedure for the linear probing diagnostic $\Phi_{\mathrm{P}}^\ell$ introduced in \cref{sec:method}.

% ------------------------------------------------------------------
\subsection{Hidden State Extraction}
\label{app:subsec:probing_extraction}

For each task instance with source prompt $S = [C;\, Q]$, we perform a single forward pass through the target model $\mathcal{M}$ and extract the hidden state $\mathbf{h}^\ell \in \mathbb{R}^d$ at the \textbf{last token position} of every layer $\ell \in \{0, 1, \dots, L{-}1\}$. Due to the causal attention mask, the last token position aggregates information from all preceding tokens, making it the natural locus for probing what the model ``knows'' at each layer.

Each hidden state $\mathbf{h}^\ell$ is paired with a label $t^* \in \mathcal{T} = \{0, 1, \dots, 9, \bar{d}\}$, where $t^*$ is the ground-truth digit and $\bar{d}$ is never used as a ground-truth label---it exists only as a prediction target (see below).

% ------------------------------------------------------------------
\subsection{11-Class Classification Setup}
\label{app:subsec:probing_classes}

A standard probing setup for this benchmark would use 10-class classification over $\mathcal{D} = \{0, 1, \dots, 9\}$. We extend this to \textbf{11 classes} by introducing the residual class $\bar{d}$, which aggregates all non-digit tokens into a single category.

\paragraph{Training label construction.} Ground-truth labels are always digit classes ($t^* \in \{0, \dots, 9\}$). To construct training samples for the $\bar{d}$ class, we include hidden states from layers where the model's final output token (via standard forward pass) is a non-digit token. This ensures the probe learns to distinguish between ``the representation points toward a digit'' and ``the representation has not yet consolidated toward any digit answer.''

\paragraph{Rationale for $\bar{d}$.} The residual class serves as an \emph{error sink}: if $\arg\max \Phi_{\mathrm{P}}^\ell(\mathbf{h}^\ell) = \bar{d}$, the probe indicates that no digit is yet linearly separable in the activation space---the information is either absent or encoded in a format that a linear classifier cannot access. This gives the probing diagnostic a three-valued interpretation at each layer: (1)~correct digit, (2)~wrong digit, (3)~not yet a digit. Without $\bar{d}$, the probe would be forced to assign probability mass to some digit even when the representation carries no digit-relevant signal, inflating false positives in early layers.

% ------------------------------------------------------------------
\subsection{Per-Layer Classifier Training}
\label{app:subsec:probing_training}

We train \textbf{one independent classifier per layer} $\ell$, denoted $\Phi_{\mathrm{P}}^\ell$. Each classifier is a logistic regression model (multinomial, L2-regularized):
%
\begin{equation}
\Phi_{\mathrm{P}}^\ell(\mathbf{h}^\ell) = \mathrm{softmax}(W^\ell \mathbf{h}^\ell + \mathbf{b}^\ell)
\end{equation}
%
where $W^\ell \in \mathbb{R}^{11 \times d}$ and $\mathbf{b}^\ell \in \mathbb{R}^{11}$.

\paragraph{Implementation.} We use \texttt{sklearn.linear\_model.LogisticRegression} with the configuration in \cref{tab:probe_hyperparams}.

\begin{table}[h]
\centering
\resizebox{\columnwidth}{!}{%
\begin{tabular}{lll}
\toprule
\textbf{Hyperparameter} & \textbf{Value} & \textbf{Rationale} \\
\midrule
Solver        & \texttt{lbfgs}       & Standard for multinomial logistic regression \\
Regularization & L2, $C = 1.0$       & Prevents overfitting on high-dim.\ $\mathbf{h}^\ell$ \\
Max iterations & 1000                & Ensures convergence across all layers \\
Multi-class   & \texttt{multinomial}  & Joint optimization over all 11 classes \\
\bottomrule
\end{tabular}%
}
\caption{Probing classifier hyperparameters.}
\label{tab:probe_hyperparams}
\end{table}

\paragraph{Why logistic regression?} The choice of a linear classifier is deliberate and methodologically motivated. A linear probe tests whether the target information is \textbf{linearly separable} in the representation space---the weakest geometric condition for information presence. Using a more powerful classifier (e.g., MLP) would conflate ``the information exists'' with ``a nonlinear transformation can extract it,'' weakening the interpretive value of the necessary condition that probing is designed to establish. This follows established practice in representation probing \citep{belinkov2022probing, hewitt2019designing}.

% ------------------------------------------------------------------
\subsection{Train/Test Split}
\label{app:subsec:probing_split}

For each task family, the dataset is split into training and test sets. The probing classifier is trained on the training split and evaluated on the held-out test split. Importantly, the \textbf{same test split} is used for both probing ($\Phi_{\mathrm{P}}^\ell$) and activation verbalization ($\Phi_{\mathrm{C}}^\ell$), ensuring that the two diagnostics are evaluated on identical instances and their layer-wise agreement/disagreement is computed on matched samples.

% ------------------------------------------------------------------
\subsection{Evaluation Metrics}
\label{app:subsec:probing_eval}

For each layer $\ell$ and each test instance, we record:
\begin{itemize}
    \item $\hat{t}_{\mathrm{P}}^\ell = \arg\max \Phi_{\mathrm{P}}^\ell(\mathbf{h}^\ell)$: the probe's predicted class.
    \item Whether $\hat{t}_{\mathrm{P}}^\ell = t^*$ (correct), $\hat{t}_{\mathrm{P}}^\ell \in \mathcal{D} \setminus \{t^*\}$ (wrong digit), or $\hat{t}_{\mathrm{P}}^\ell = \bar{d}$ (non-digit).
\end{itemize}

The \textbf{First Probing Correct Layer (FPCL)} is defined as:
%
\begin{equation}
\mathrm{FPCL}(S) = \min \{ \ell \mid \hat{t}_{\mathrm{P}}^\ell = t^* \}
\end{equation}

The \textbf{Brewing Duration} $\Delta$ is the gap between FPCL and the First Correct Layer of CSD:
%
\begin{equation}
\Delta(S) = \mathrm{FCL}(S) - \mathrm{FPCL}(S)
\end{equation}

$\Delta > 0$ indicates that the information exists in the representation (the probe can read it) before the model's own decoding pathway can access it---the hallmark of \textbf{Information Brewing}.
